% ===========================================================================
%  PART 2 --- Rotations and Insertion
%  Flow: one rotation drawn as a physical turn -> the four cases table ->
%        the 2x2 grid of pivots -> the insertion story (which ends on a
%        question) -> one simulated insert that answers it -> the catch ->
%        why one rotation is always enough, and what it costs
%
%  Rotation language, used identically on every rotation in this deck:
%    a PIN goes into the node the rotation is NAMED at (the broken one), and
%    one big ARC over the tree shows which way the whole thing turns --
%    right rotation = clockwise, left rotation = anticlockwise.
%  The result is always a SEPARATE overlay, with ghosts at the old positions.
% ===========================================================================

\avldivider{2}{Rotations}{How the tree fixes itself.}{avlptwo}{avlptwoA}{avlptwoB}{avlptwoC}

\setavlpart{2}{Part 2 --- Rotations}{avlptwo}

% ------------------------------------------------------------------ ONE ROTATION
%  The whole idea in one frame: pin a node, turn the tree about it.
%
%  Layout rule for this frame: a node's x is its position in sorted order and
%  NEVER changes; only its depth does. A rotation therefore moves nodes
%  straight up or straight down, which is what makes the ghosts readable --
%  every faded circle sits directly above or below where its node landed. It
%  also makes the payoff visible rather than asserted: nothing moves sideways,
%  so the search order cannot break.
\begin{frame}[label=p2]{A rotation turns the tree about a pin}
  \vspace{-0.2em}
  \begin{center}
  \begin{tikzpicture}[x=1cm,y=1cm,scale=0.95,every node/.style={transform shape}]
    % fixed extent: the picture must not jump between overlay steps. The top is
    % set by the arc (radius 1.65 about a node at y = 0.8).
    \path (-3.9,-2.78) rectangle (3.9,2.55);

    % ---------------- before: 30, 20, 10 all went left ----------------
    \only<1-2>{
      \draw[avledge] (1.9,0.8) -- (0,-0.5) -- (-1.9,-1.8);
      \node[avlkeyalert] at (1.9,0.8)   {30};
      \node[avlkey]      at (0,-0.5)    {20};
      \node[avlkeygold]  at (-1.9,-1.8) {10};
    }
    \only<1>{%
      \draw[avlalert,line width=2\pxl] (1.9,0.8) circle (0.42);
      \node[avlannotalert,anchor=west] at (2.42,0.98) {$\mathrm{bf} = +2$};}
    % Pin and turn are ONE beat: the pin is only meaningful as the thing the
    % arc sweeps about, and splitting them made the audience click twice for
    % half a gesture.
    \only<2>{\pivotpin{1.9,0.8}\rotright{1.9,0.8}{1.65}}

    % ---------------- after: 20 came up, 30 came down -----------------
    \only<3->{
      % Afterimages on the result beat ONLY. They carry the eye across the jump
      % cut from the turn; by the next beat the order strip is making the point
      % instead, and three dashed discs would just read as extra nodes.
      \only<3>{%
        \ghostnode{1.9}{0.8}{30}  \ghostpath{1.9,0.45}{1.9,-0.15}
        \ghostnode{0}{-0.5}{20}   \ghostpath{0,-0.15}{0,0.45}
        \ghostnode{-1.9}{-1.8}{10} \ghostpath{-1.9,-1.45}{-1.9,-0.85}}
      % restored: the tree that comes out of the turn is legal, so it is green
      \draw[draw=avlbalance!70,line width=1.6\pxl] (-1.9,-0.5) -- (0,0.8) -- (1.9,-0.5);
      \node[avlkeyok] at (0,0.8)     {20};
      \node[avlkeyok] at (-1.9,-0.5) {10};
      \node[avlkeyok] at (1.9,-0.5)  {30};
      % No pin on these beats. The pin was a fixed post; redrawing it under 30
      % after the turn would say it travelled with the node, which is the one
      % thing the whole picture is trying not to say.
    }

    % ---------------- the order strip, from the payoff beat on --------
    \only<4->{
      \draw[avlrule,line width=1\pxl] (-2.6,-2.20) -- (2.6,-2.20);
      \foreach \k/\x in {10/-1.9, 20/0, 30/1.9} {
        \draw[avlrule,line width=1\pxl] (\x,-2.08) -- (\x,-2.32);
        \node[color=avlink,font=\avlfAnnot] at (\x,-1.90) {\k};
      }
      \node[avlannot,anchor=north] at (0,-2.42)
        {left to right order --- \strong{never changed}};
    }
  \end{tikzpicture}
  \end{center}

  \vspace{-0.3em}
  \begin{center}
    \only<1>{{\avlfSec\color{avlmuted}Left skewed tree.
              \viol{30 has a balance factor of $+2$}.}}
    \only<2>{{\avlfSec\color{avlmuted}Push the pin into
              {\color{avlptwo}\wmed 30} and turn the whole tree
              {\color{avlptwo}\wmed clockwise} about it. That is a
              {\color{avlptwo}\wmed right rotation}.}}
    \only<3>{{\avlfSec\color{avlmuted}20 became the parent, 30 came down on the
              right. \ok{The tree is balanced again.}}}
    \only<4->{{\avlfSec\color{avlmuted}Nothing moved sideways, so the order is
              safe, and only \fib{three links} changed.}}
  \end{center}

  \note{Say it as a physical action. The pin goes into 30, the node that is
  broken, and in the same breath you turn the whole tree clockwise about that
  pin, the way you would turn a stiff wire. One gesture, not two.
  20 ends up on top, 30 ends up on the right. Point at
  the faded circles: that is where each node was a second ago. Then point at
  the ruler: nothing slid left or right. That is why a rotation can never break
  the search rule, and it only ever re-links three pointers, so it is constant
  time.}
\end{frame}
% ------------------------------------------------------------------ four cases
\begin{frame}[label=p2]{Four cases, and what to do}
  % No vertical rules, no zebra. Header labels avlfaint, hairline under the
  % header, avlrule2 between body rows, the fix column carries the ember tint.
  % One hairline under the header and nothing else: no vertical rules, no row
  % rules, no fills. The fix column is carried by colour alone.
  \vspace{0.1em}
  \begin{center}
  \setlength{\tabcolsep}{20\pxl}
  \renewcommand{\arraystretch}{1.32}
  \setlength{\arrayrulewidth}{1\pxl}
  \arrayrulecolor{avlrule}
  {\avlfBody
  \begin{tabular}{lll}
    {\avlfEyebrow\color{avlfaint}case}
      & {\avlfEyebrow\color{avlfaint}where the new node went}
      & {\avlfEyebrow\color{avlfaint}fix}\\\hline
    \strong{LL} & left child's \emph{left} side  & \fib{one right rotation}\\
    \strong{RR} & right child's \emph{right} side & \fib{one left rotation}\\
    \strong{LR} & left child's \emph{right} side & \fib{left, then right}\\
    \strong{RL} & right child's \emph{left} side & \fib{right, then left}\\
  \end{tabular}}
  \end{center}

  \vspace{0.5em}
  \begin{center}
    {\avlfSec\color{avlmuted}Each case is named after the two steps the new
     node took down from the broken node.}\\[0.7em]
    \begin{avlpanel}[0.96\textwidth]
      A \fib{zig-zag} imbalance needs two rotations: first on the child, then
      on the broken node.
    \end{avlpanel}
  \end{center}

  \note{Four shapes, four fixes. The two straight ones need a single rotation.
  The two bent ones need two rotations: first straighten the bend, then rotate
  as normal. Do not linger --- the next slide shows all four side by side.}
\end{frame}

% ------------------------------------------------------------------ 2x2 GRID
%  The four cases as four actual pivots, laid out 2x2. Same three keys every
%  time (10, 20, 30) so the only variable is arrival ORDER, and every cell ends
%  on the same balanced tree.
%
%  Cell shapes are written once here: the four cases are just the same three
%  nodes in four arrangements, and the doubles turn into the singles after
%  their first rotation, which is the point the frame is making.
% Cell lattice, the same rule as the simulation frame: 10 always sits at
% x = -0.90, 20 at x = 0, 30 at x = +0.90, in every shape. A click therefore
% moves a node straight up or straight down and never sideways, and the bent
% shapes visibly straighten into the chains directly above them.
\newcommand{\cellchainL}{%   leans left -- the LL shape
  \draw[avledge,line width=0.9pt] (0.90,0) -- (0,-0.80) -- (-0.90,-1.60);
  \node[gred] at (0.90,0) {30}; \node[gkey] at (0,-0.80) {20};
  \node[ggold] at (-0.90,-1.60) {10};}
\newcommand{\cellchainR}{%   leans right -- the RR shape
  \draw[avledge,line width=0.9pt] (-0.90,0) -- (0,-0.80) -- (0.90,-1.60);
  \node[gred] at (-0.90,0) {10}; \node[gkey] at (0,-0.80) {20};
  \node[ggold] at (0.90,-1.60) {30};}
% The two bent shapes do NOT use the chain's even level drop. A bend spans two
% lattice columns on its upper edge and only one on its lower edge, so at an
% even 0.80/0.80 drop the upper edge comes out 1.97 long against a 1.20 stub --
% which is what made LR and RL read lopsided next to the chains. The column
% positions are fixed by sorted order and cannot move, so the only free lever
% is the drop: 0.52 then 1.08 keeps the total height at 1.60 (the tags and the
% row hairline never move) and brings the two edges to 1.87 against 1.42.
% The bend node is the ONLY node on its level in these shapes, so an uneven
% drop is invisible -- there is nothing beside it to measure against.
\newcommand{\cellbentL}{%    left then right -- the LR shape
  \draw[avledge,line width=0.9pt] (0.90,0) -- (-0.90,-0.52) -- (0,-1.60);
  \node[gred] at (0.90,0) {30}; \node[gkey] at (-0.90,-0.52) {10};
  \node[ggold] at (0,-1.60) {20};}
\newcommand{\cellbentR}{%    right then left -- the RL shape
  \draw[avledge,line width=0.9pt] (-0.90,0) -- (0.90,-0.52) -- (0,-1.60);
  \node[gred] at (-0.90,0) {10}; \node[gkey] at (0.90,-0.52) {30};
  \node[ggold] at (0,-1.60) {20};}
\newcommand{\cellflat}{%     the balanced result every case lands on
  \draw[avlgoldedge,line width=1.2pt] (-0.90,-0.80) -- (0,0) -- (0.90,-0.80);
  \node[ggold] at (0,0) {20}; \node[ggold] at (-0.90,-0.80) {10};
  \node[ggold] at (0.90,-0.80) {30};}

\begin{frame}[label=p2]{\only<1>{Four cases, four pivots}%
    \only<2>{LR and RL just become LL and RR}%
    \only<3->{Four arrivals, one tree}}
  \vspace{-1.1em}
  \begin{center}
  % The grid used to sit at scale 0.70 inside a 14 cm text block -- barely 8 cm
  % wide, with the discs and their tags shrunk to match. The lattice is now
  % drawn larger and spread wider, so the figure fills the column it is given.
  % The grid used to sit at scale 0.70 inside a 14 cm text block -- barely 8 cm
  % wide, with the discs and their tags shrunk to match. The lattice is drawn
  % larger and spread wider now, and each tag is ONE line: two-line tags were
  % what made the block too tall to enlarge without pushing into the panel.
  \begin{tikzpicture}[x=1cm,y=1cm,scale=0.72,every node/.style={transform shape},
      gkey/.style ={avlkey,      minimum size=6.6mm,font=\avlfEyebrow,inner sep=0pt},
      ggold/.style={avlkeygold,  minimum size=6.6mm,font=\avlfEyebrow,inner sep=0pt},
      gred/.style ={avlkeyalert, minimum size=6.6mm,font=\avlfEyebrow,inner sep=0pt},
      tag/.style  ={font=\avlfBody,align=center,text=avlmuted,anchor=north,
                    inner sep=1pt}]
    % hairline between the two rows of cells, clear of the row-1 tags above it
    % and the row-2 discs below it
    \only<3->{\draw[avlrule,line width=1\pxl] (-6.6,-2.68) -- (6.6,-2.68);}
    \setpinscale{1.0}
    % Fixed extent, symmetric about x = 0. Row pitch 3.32: an arc reaches 0.85
    % above its cell origin and a one-line tag reaches 2.34 below it, so
    % anything tighter puts the top row's caption inside the bottom row's arc.
    \path (-7.1,-5.66) rectangle (7.1,0.88);

    % ================= LL : 30, 20, 10 =================
    \begin{scope}[shift={(-5.2,0)}]
      % arc first, so the nodes always sit on top of it
      \only<1>{\rotright{0.90,0}{0.85}\cellchainL\pivotpin{0.90,0}}
      \only<2->{\cellflat}
      \node[tag] at (0,-2.02) {\textbf{\color{avlink}LL}\ \ 30, 20, 10
        \ \quietc{---}\ \fib{right 30}};
    \end{scope}

    % ================= RR : 10, 20, 30 =================
    \begin{scope}[shift={(5.2,0)}]
      \only<1>{\rotleft{-0.90,0}{0.85}\cellchainR\pivotpin{-0.90,0}}
      \only<2->{\cellflat}
      \node[tag] at (0,-2.02) {\textbf{\color{avlink}RR}\ \ 10, 20, 30
        \ \quietc{---}\ \fib{left 10}};
    \end{scope}

    % ================= LR : 30, 10, 20 =================
    % The first turn is at 10 and straightens the bend into an LL chain; the
    % second turn is then literally the LL cell sitting above it.
    \begin{scope}[shift={(-5.2,-3.32)}]
      \only<1>{\rotleft{-0.90,-0.52}{0.70}
               \cellbentL\pivotpin{-0.90,-0.52}}
      \only<2>{\rotright{0.90,0}{0.85}\cellchainL\pivotpin{0.90,0}}
      \only<3->{\cellflat}
      \node[tag] at (0,-2.02) {\textbf{\color{avlink}LR}\ \ 30, 10, 20
        \ \quietc{---}\ \fib{left, right}};
    \end{scope}

    % ================= RL : 10, 30, 20 =================
    \begin{scope}[shift={(5.2,-3.32)}]
      \only<1>{\rotright{0.90,-0.52}{0.70}
               \cellbentR\pivotpin{0.90,-0.52}}
      \only<2>{\rotleft{-0.90,0}{0.85}\cellchainR\pivotpin{-0.90,0}}
      \only<3->{\cellflat}
      \node[tag] at (0,-2.02) {\textbf{\color{avlink}RL}\ \ 10, 30, 20
        \ \quietc{---}\ \fib{right, left}};
    \end{scope}
  \end{tikzpicture}
  \end{center}

  \only<3->{
  \vspace{0.45em}
  \begin{center}
    \begin{avlpanel}[0.96\textwidth]
      Same three keys, four arrival orders --- \fib{all four land on the same
      tree}.
    \end{avlpanel}
  \end{center}
  \vspace{0.15em}}

  \note{Four cells, same three keys, only the arrival order changes. Point at
  each pin and each arc: LL turns clockwise about 30, RR turns anticlockwise
  about 10. Click. LL and RR are done. Now the important beat: the bent ones
  are no longer bent --- LR has turned into an LL chain and RL into an RR
  chain. Click again. Same pin, same turn as the cell directly above, and all
  four land on the same tree.}
\end{frame}

% ------------------------------------------------------------------ rotation power
\begin{frame}[label=p2]{The beauty: one constant operation fixes everything}
  \vspace{-0.4em}
  \begin{center}
  \eyebrow{What makes a rotation elegant}
  \end{center}

  \vspace{0.5em}
  \begin{center}
  \begin{minipage}{0.86\textwidth}
  \begin{columns}[T,onlytextwidth]
    \begin{column}{0.47\textwidth}
      {\fib{Only three links change}}\\[0.15em]
      {\avlfSec\color{avlmuted}$O(1)$ no matter the tree size}\\[0.8em]
      {\fib{Restores subtree height}}\\[0.15em]
      {\avlfSec\color{avlmuted}Parent sees no change, repair stops}
    \end{column}
    \begin{column}{0.47\textwidth}
      {\fib{Search order is safe}}\\[0.15em]
      {\avlfSec\color{avlmuted}Left-right order never breaks}\\[0.8em]
      {\fib{One rotation suffices}}\\[0.15em]
      {\avlfSec\color{avlmuted}No need to check nodes above}
    \end{column}
  \end{columns}
  \end{minipage}
  \end{center}

  \vspace{1.0em}
  \begin{center}
    \begin{avlpanel}[0.84\textwidth]
      Three pointers, no recursion, no loops --- that is why a rotation
      \fib{runs in constant time}.
    \end{avlpanel}\\[0.7em]
    {\avlfSec\color{avlmuted}Four cases. \quad One pattern. \quad
     \ok{Perfect balance.}}
  \end{center}

  \note{The heart of AVL. A rotation is O(1): only three pointers move. It
  restores the subtree to its original height, so the parent tree sees nothing
  has changed. That is why one rotation is always enough for insertion. This is
  the elegant core that makes AVL practical.}
\end{frame}

% ------------------------------------------------------------------ insertion story
%  The insertion arc, four frames:
%    1. the story        insert -> walk up -> rotate, ending on the QUESTION
%    2. the simulation   one insert on a tree nobody has seen, answering it
%    3. the catch        one rotation? always? yes
%    4. why so           the height argument, then the cost
%
%  The question is asked before it is answered on purpose. The old deck
%  asserted "one rotation is always enough" on the rule slide and then ran a
%  simulation that could only confirm what had already been given away.
%
%  WORDING, fixed across all three frames up to the catch: the third step is
%  "rotate when we find a bad node". Not "the first bad node" -- that phrasing
%  quietly promises there is only one, which is the thing the catch is about
%  to ask.
\begin{frame}[label=p2]{Insertion: insert, walk up, rotate}
  \vspace{0.7em}
  \begin{center}
  \begin{tikzpicture}[x=1cm,y=1cm,
      stepbox/.style={draw=avlptwo!45!avlground,fill=avlptwo!8!avlground,
                   line width=1.5\pxl,rounded corners=6pt,
                   minimum width=3.4cm,minimum height=1.15cm,
                   font=\avlfSec,text=avlink},
      last/.style={draw=avlfib!45!avlground,fill=avlfib!8!avlground,
                   line width=1.5\pxl,rounded corners=6pt,
                   minimum width=3.4cm,minimum height=1.15cm,
                   font=\avlfSec,text=avlink}]
    \node[stepbox] (s1) at (-4.4,0)
      {\begin{minipage}{3.0cm}\centering Insert like a normal BST\end{minipage}};
    \node[stepbox] (s2) at (0,0)
      {\begin{minipage}{3.0cm}\centering Walk back up to the root\end{minipage}};
    \node[last] (s3) at (4.4,0)
      {\begin{minipage}{3.0cm}\centering Rotate when we find a bad node\end{minipage}};
    \draw[avlgoldarrow] (s1) -- (s2);
    \draw[avlgoldarrow] (s2) -- (s3);
  \end{tikzpicture}
  \end{center}

  \vspace{0.9em}
  \begin{center}
  % Fixed-height box so the flow diagram above does not shift when the
  % question replaces the caption.
  \begin{minipage}[t][150\pxl][t]{0.96\textwidth}
  \centering
    \only<1>{{\avlfSec\color{avlmuted}The tree is only repaired on the way
              \emph{back up}.}\\[0.9em]
             {\avlfSec\color{avlmuted}A node is bad once its balance factor
              reaches $\pm 2$.}}
    \only<2->{%
      \begin{avlpanel}[0.78\textwidth]
        {\avlfLead\color{avlfib}\wbold But exactly how many rotations?}
      \end{avlpanel}\\[0.8em]
      {\avlfSec\color{avlmuted}Let us run one insert and count.}}
  \end{minipage}
  \end{center}

  \note{Three steps, and the order matters: you insert like an ordinary BST,
  then you walk back up, and when you find a node with a balance factor of plus
  or minus two, you rotate there. Click. Now hold the question open: how many
  rotations does that cost? Do not answer it yet, the next slide answers it by
  counting.}
\end{frame}
% ------------------------------------------------------------------ INSERTION SIM
%  ONE insert into a tree nobody has seen before, so the count is earned rather
%  than asserted. The tree is chosen so the walk up passes two healthy nodes
%  before it finds the broken one: the whole point of step 2 is that most of the
%  walk finds nothing, and a tree that broke at the parent would hide that.
%
%    start        40(20(10,30),60)          legal: bf(40)=+1, bf(20)=0
%    insert 5     lands left of 10          bf(10)=+1 ok, bf(20)=+1 ok,
%                                           bf(40)=+2 BROKEN, an LL shape
%    fix          one right rotation at 40  -> 20(10(5),40(30,60))
%
%  Height check, which is the reason one rotation ends it: the tree stood at
%  height 2 before the insert and stands at height 2 after the rotation.
%
%  LAYOUT: lattice pitch 1.25 x 1.20 at scale 0.84, the same figures the 1..7
%  simulation used, so this tree reads at the size the deck has already
%  established. A key's COLUMN is its position in sorted order and never
%  changes, only its ROW does. c0..c5 = 5, 10, 20, 30, 40, 60.
%
%  There is no caption strip under this figure. Two things were competing for
%  the same 1 cm: a readable tree and a line of prose. The tree won, and what
%  the prose was carrying (the balance factors, the shape name, the count) is
%  annotated on the figure itself where it belongs.
%
%  ALGORITHM RAIL. All four steps hold their slots from beat one, so the list
%  is a map of where we are rather than a list being built. \setbeamercovered
%  does the dimming: a step outside its own overlays is covered, which beamer
%  renders at 32% opacity, and \alt paints it muted grey underneath. The green
%  chevron is drawn only on the active beats. "Done" is the exception: it is
%  \only-gated, so it is absent, not dimmed, until the tree is actually legal.

% \insmark{<active spec>} -- the "you are here" chevron. Always occupies its
% slot (the \path reserves the box) and is only inked on the active beats, so
% the rail cannot shift sideways between overlays. 1 unit = 1 design px.
\newcommand{\insmark}[1]{%
  \begin{tikzpicture}[baseline=0pt,x=0.0125cm,y=0.0125cm]
    \path (0,0) rectangle (11,12);
    \only<#1>{\fill[avlbalance] (0,1) -- (0,11) -- (9.5,6) -- cycle;}
  \end{tikzpicture}}

% \insstep[<reveal spec>]{<active spec>}{<text>}
\newcommand{\insstep}[3][1-]{%
  \hbox to\linewidth{%
    \insmark{#2}\hskip14\pxl
    \begin{minipage}[t]{\dimexpr\linewidth-40\pxl\relax}
      \avlfSec\raggedright\strut
      \only<#1>{\uncover<#2>{%
        \alt<#2>{\color{avlbalance}\wsemi}{\color{avlmuted}}#3}}%
    \end{minipage}\hfil}%
  \vskip 24\pxl}

\begin{frame}[label=p2]{Simulation: count the rotations}
  \setbeamercovered{transparent=32}
  \vspace{0.5em}
  \begin{columns}[T,onlytextwidth]
  \begin{column}{0.60\textwidth}
  \begin{center}
  \begin{tikzpicture}[x=1cm,y=1cm,scale=0.84,every node/.style={transform shape}]
    \setpinscale{1}
    % column = sorted position (5,10,20,30,40,60), row = depth
    \foreach \cc in {0,...,5}{\foreach \dd in {0,1,2,3}{
      \coordinate (c\cc-\dd) at ({(\cc-2.5)*1.25},{-\dd*1.20});}}

    % Fixed extent: the tree must not resize between beats. The top is set by
    % the turn arrow and its label, the bottom by the ghost of 5.
    \path (-3.70,-3.98) rectangle (3.70,1.50);

    % ---------------- the tree as it stands once 5 has arrived -------------
    \only<1-4>{
      \draw[avledge] (c4-0) -- (c2-1) -- (c1-2) -- (c0-3);
      \draw[avledge] (c4-0) -- (c5-1);
      \draw[avledge] (c2-1) -- (c3-2);
    }
    % the walk itself, drawn on the links it climbs
    \only<2-4>{\draw[draw=avlptwo,line width=2\pxl] (c0-3) -- (c1-2) -- (c2-1);}
    \only<3-4>{\draw[draw=avlptwo,line width=2\pxl] (c2-1) -- (c4-0);}

    \only<1>{
      \node[avlkey]     at (c4-0) {40}; \node[avlkey] at (c2-1) {20};
      \node[avlkey]     at (c5-1) {60}; \node[avlkey] at (c1-2) {10};
      \node[avlkey]     at (c3-2) {30}; \node[avlkeygold] at (c0-3) {5};
      \node[avlannotfib,anchor=west] at (-2.62,-3.60) {new};
    }
    % Walk beats. A node that has been checked and passed keeps its green, so
    % the audience can see how much of the walk found nothing.
    \only<2-4>{
      \node[avlkeyok] at (c2-1) {20}; \node[avlkey] at (c5-1) {60};
      \node[avlkeyok] at (c1-2) {10}; \node[avlkey] at (c3-2) {30};
      \node[avlkeygold] at (c0-3) {5};
      \node[avlannotok,anchor=east] at (-2.38,-2.40) {$\mathrm{bf} = +1$};
      \node[avlannotok,anchor=east] at (-1.13,-1.20) {$\mathrm{bf} = +1$};
    }
    \only<2>{\node[avlkey] at (c4-0) {40};}
    \only<3-4>{\node[avlkeyalert] at (c4-0) {40};}
    \only<3>{
      \draw[avlalert,line width=2\pxl] (c4-0) circle (0.42);
      \node[avlannotalert,anchor=west] at (2.45,0.18) {$\mathrm{bf} = +2$};}
    \only<4>{
      \pivotpin{c4-0}\rotright{c4-0}{1.30}
      \node[color=avlptwo,font=\avlfAnnot\wmed,anchor=east] at (0.40,1.28)
        {LL: turn right at 40};}

    % ---------------- the turn has landed, afterimages still up ------------
    \only<5>{
      \ghostnode{1.875}{0}{40}       \ghostpath{1.875,-0.35}{1.875,-0.85}
      \ghostnode{-0.625}{-1.20}{20}  \ghostpath{-0.625,-0.85}{-0.625,-0.35}
      \ghostnode{-3.125}{-3.60}{5}   \ghostpath{-3.125,-3.25}{-3.125,-2.75}
      \draw[avlgoldedge,line width=1.1pt] (c1-1) -- (c2-0) -- (c4-1);
      \draw[avlgoldedge,line width=1.1pt] (c0-2) -- (c1-1);
      \draw[avlgoldedge,line width=1.1pt] (c3-2) -- (c4-1) -- (c5-2);
      \node[avlkeygold] at (c2-0) {20}; \node[avlkeygold] at (c1-1) {10};
      \node[avlkeygold] at (c4-1) {40}; \node[avlkeygold] at (c0-2) {5};
      \node[avlkeygold] at (c3-2) {30}; \node[avlkeygold] at (c5-2) {60};
    }
    % ---------------- verdict: same tree, afterimages cleared, green -------
    % Nothing moves between beat 5 and beat 6. The only change is that the
    % tree stops being a thing in motion and becomes a legal result.
    \only<6>{
      \draw[draw=avlbalance!70,line width=1.6\pxl] (c1-1) -- (c2-0) -- (c4-1);
      \draw[draw=avlbalance!70,line width=1.6\pxl] (c0-2) -- (c1-1);
      \draw[draw=avlbalance!70,line width=1.6\pxl] (c3-2) -- (c4-1) -- (c5-2);
      \node[avlkeyok] at (c2-0) {20}; \node[avlkeyok] at (c1-1) {10};
      \node[avlkeyok] at (c4-1) {40}; \node[avlkeyok] at (c0-2) {5};
      \node[avlkeyok] at (c3-2) {30}; \node[avlkeyok] at (c5-2) {60};
      \node[font=\avlfLead\wbold,text=avlbalance,anchor=north] at (0,-3.05)
        {Legal again. One rotation.};
    }
  \end{tikzpicture}
  \end{center}
  \end{column}

  % ---------------- the algorithm, tracked beside the tree ----------------
  \begin{column}{0.36\textwidth}
  \vspace{1.1em}
  \insstep{1}{Insert like a normal BST}
  \insstep{2-3}{Walk back up to the root}
  \insstep{4-5}{Rotate when we find a bad node}
  \insstep[6]{6}{Done}
  \end{column}
  \end{columns}

  \note{A tree they have not seen. Click: 5 arrives, ordinary BST insert, no
  balance work at all yet. Click: now walk back up. 10 is fine, 20 is fine.
  Say out loud that this is the common case, the walk usually finds nothing.
  Click: the root is the one that broke, plus two. Click: left child, left
  side, so it is LL. Pin 40, turn clockwise. Click: the turn has landed, and
  the faded circles say where 40, 20 and 5 came from. Click: legal. Count it
  with them, one rotation. Then ask whether that was luck.}
\end{frame}

% ------------------------------------------------------------------ the catch
%  ONE box on this slide, and it is the answer. The questions are set as plain
%  type with real air around them, so the only bordered thing on the ground is
%  the thing the audience is waiting for. The hierarchy is carried by size --
%  17 pt questions against a 25 pt answer -- and by colour: ink for the
%  question anyone would ask, part hue for the one that presses, green for the
%  verdict. A second card around the questions would have flattened all of it
%  back into two grey rectangles.
\begin{frame}[label=p2]{The catch}
  \vspace{0.6em}
  \begin{center}
    % The baselineskip in force here is \avlfBody's 14.4 pt, not \avlfStmt's:
    % the size command is scoped inside the braces, so the \\ that follows it
    % still breaks at body leading. That is why these gaps look outsized.
    {\avlfStmt\wmed\color{avlink}One rotation and we are done?}\\[72\pxl]
    \uncover<2->{{\avlfStmt\wmed\color{avlptwo}Always the case?}}\\[78\pxl]
    % Fixed slot: the answer must not push the questions up the slide when it
    % arrives. Sized to the panel plus its trailing line.
    \begin{minipage}[t][174\pxl][t]{\textwidth}
    \centering
    \only<3->{%
      \begin{avlpanelh}[0.46\textwidth]{avlbalance}
        {\avlfBig\wbold\color{avlbalance}Yes, it is.}
      \end{avlpanelh}\\[18\pxl]
      {\avlfSec\color{avlmuted}At most one rotation per insert.}}
    \end{minipage}
  \end{center}

  \note{Slow down here. One rotation. Was that this tree, or is it a
  guarantee? Let them sit with it for a beat. Click. Ask it again, out loud:
  always? Click. Yes, always, at most one rotation per insert, and that is not
  an observation, it is provable. The next slide is the proof in one picture.}
\end{frame}

% ------------------------------------------------------------------ why so
%  The height argument, drawn rather than asserted: the insert lifts one
%  subtree by exactly one, the rotation puts that one back, and an ancestor
%  that sees the height it saw before has nothing left to repair.
\begin{frame}[label=p2]{Why so?}
  \vspace{-0.1em}
  \begin{center}
  \begin{tikzpicture}[x=1cm,y=1cm,scale=0.75,every node/.style={transform shape},
      par/.style={circle,draw=avlink!70,fill=white!96!avlground,
                  line width=1.5\pxl,inner sep=0pt,minimum size=30\pxl},
      whycap/.style={font=\avlfAnnot,text=avlmuted,anchor=north,align=center}]
    \path (-6.30,-1.80) rectangle (6.30,1.45);

    % ---- before the insert : height h ----
    \begin{scope}[shift={(-4.60,0)}]
      \node[par] at (0,1.20) {};
      \draw[avledge] (0,1.20) -- (0,0.85);
      \draw[draw=avlfib!55,fill=avlfib!8!avlground,line width=1.6\pxl]
        (0,0.85) -- (-1.00,-0.85) -- (1.00,-0.85) -- cycle;
      \heightbrace{-1.32}{0.85}{-0.85}{$h$}{avlfaint}
      \node[whycap] at (0,-1.05) {before the insert};
    \end{scope}

    % ---- after the insert : one taller ----
    \begin{scope}
      \node[par] at (0,1.20) {};
      \draw[avledge] (0,1.20) -- (0,0.85);
      \draw[draw=avlalert!55,fill=avlalert!8!avlground,line width=1.6\pxl]
        (0,0.85) -- (-1.15,-1.25) -- (1.15,-1.25) -- cycle;
      \heightbrace{-1.47}{0.85}{-1.25}{$h+1$}{avlalert}
      \node[whycap] at (0,-1.45) {after the insert};
    \end{scope}

    % ---- after one rotation : height h again ----
    % Pushed a little further right than the mirror position: this brace sits
    % on the same y as the "rotate" arrow, and its rotated label needs clear
    % air between itself and the arrowhead.
    \begin{scope}[shift={(4.80,0)}]
      % Whole state hidden until the rotation beat -- the parent stub on its
      % own reads as a stray lollipop with no subtree under it.
      \only<2->{
        \node[par] at (0,1.20) {};
        \draw[avledge] (0,1.20) -- (0,0.85);
        \draw[draw=avlbalance!60,fill=avlbalance!8!avlground,line width=1.6\pxl]
          (0,0.85) -- (-1.00,-0.85) -- (1.00,-0.85) -- cycle;
        \heightbrace{-1.32}{0.85}{-0.85}{$h$}{avlbalance}
        \node[whycap] at (0,-1.05) {after one rotation};}
    \end{scope}

    \draw[avlgoldarrow] (-3.35,-0.10) -- (-1.95,-0.10)
      node[midway,above=2\pxl,font=\avlfAnnot,text=avlfib] {insert};
    \only<2->{
      \draw[avlarc] (1.70,-0.10) -- (2.90,-0.10)
        node[midway,above=2\pxl,font=\avlfAnnot,text=avlptwo] {rotate};}
  \end{tikzpicture}
  \end{center}

  \vspace{0.2em}
  \begin{center}
  % HARD BUDGET. Below the figure there are 2.45 cm, and a fixed-height
  % minipage hides an overrun from the compiler: a wrapped line here spills
  % silently through the footer rule instead of raising an overfull warning.
  % Both sentences and the panel line must each stay on ONE line, which is
  % roughly 72 characters at \avlfSec and roughly 46 at \avlfLead.
  \begin{minipage}[t][196\pxl][t]{0.94\textwidth}
  \centering
    {\avlfSec\color{avlmuted}An insert adds
     \viol{at most one level} to a subtree.}\\[0.5em]
    \uncover<2->{{\avlfSec\color{avlmuted}The rotation takes that level back,
      so \ok{nothing above the node changes}.}}\\[0.6em]
    \only<3->{%
      \begin{avlpanel}[0.94\linewidth]
        {\avlfLead\color{avlfib}\wbold Insertion is $O(\log n)$: a walk and
         one turn.}
      \end{avlpanel}}
  \end{minipage}
  \end{center}

  \note{Here is the whole argument in one picture. An insert adds one node, so
  it can lift a subtree by at most one level. Click. The rotation gives that
  level straight back: the subtree comes out at exactly the height it had
  before the insert, so the parent, and everything above the parent, sees no
  change at all and has nothing to fix. That is why the repair stops at the
  node we rotated. Click. And the cost: log n down, log n back up, one
  constant-time rotation. Insertion is logarithmic, guaranteed.}
\end{frame}
